\section{Research problem}
\label{sec:conclusions_research_problem}

Over the past decade, \acp{bss} have become an increasingly important component of urban mobility as more systems are implemented and existing ones expand, making them a relevant part of how people move around cities. This growth has been accompanied by the deployment of digital infrastructure that produces rich operational data, creating new opportunities to improve both long-term system design and day-to-day operations through computational, data-driven decision support.

Considerable progress has been made in integrating data and computational methods into \ac{bss} planning and operations, including work on system design, demand modelling, and rebalancing. Nevertheless, important gaps remain. Building on prior advances, this thesis asks how \ac{bss} data can be most effectively harnessed to improve efficiency and sustainability while supporting policy-relevant objectives, and answers it by developing computational methodologies in three areas:

\begin{enumerate}
    \item Developing a single flexible station-planning framework that can explore multiple, competing planning objectives and is detailed enough to support practical, street-level system design.
    \item Predicting \ac{e-b} battery levels without exact routes or in-transit measurements, to improve bike availability and operational efficiency.
    \item Enabling proactive and interpretable fleet maintenance, contributing to operational efficiency and more sustainable service provision.
\end{enumerate}